\begin{lemma}
\label{lemma-pseudo-coherent-affine}
Let $X = \Spec(A)$ be an affine scheme. Let $M^\bullet$ be a
complex of $A$-modules and let $E$ be the corresponding object
of $D(\mathcal{O}_X)$. Then $E$ is an $m$-pseudo-coherent
(resp.\ pseudo-coherent) as an object of $D(\mathcal{O}_X)$
if and only if $M^\bullet$ is $m$-pseudo-coherent (resp.\ pseudo-coherent)
as a complex of $A$-modules.
\end{lemma}

\begin{proof}
It is immediate from the definitions that if $M^\bullet$ is
$m$-pseudo-coherent, so is $E$. To prove the converse, assume
$E$ is $m$-pseudo-coherent. As $X = \Spec(A)$ is quasi-compact with
a basis for the topology given by standard opens, we can find a standard
open covering $X = D(f_1) \cup \ldots \cup D(f_n)$ and strictly
perfect complexes $\mathcal{E}_i^\bullet$ on $D(f_i)$ and
maps $\alpha_i : \mathcal{E}_i^\bullet \to E|_{U_i}$ inducing
isomorphisms on $H^j$ for $j > m$ and surjections on $H^m$.
By Cohomology, Lemma \ref{cohomology-lemma-local-actual}
after refining the open covering
we may assume $\alpha_i$ is given by a map of complexes
$\mathcal{E}_i^\bullet \to \widetilde{M^\bullet}|_{U_i}$
for each $i$. By Modules, Lemma
\ref{modules-lemma-direct-summand-of-locally-free-is-locally-free}
the terms $\mathcal{E}_i^n$ are finite locally free modules.
Hence after refining the open covering we may assume each
$\mathcal{E}_i^n$ is a finite free $\mathcal{O}_{U_i}$-module.
From the definition it follows that $M^\bullet_{f_i}$ is
an $m$-pseudo-coherent complex of $A_{f_i}$-modules.
We conclude by applying
More on Algebra, Lemma \ref{more-algebra-lemma-glue-pseudo-coherent}.

\medskip\noindent
The case ``pseudo-coherent'' follows from the fact that $E$ is
pseudo-coherent if and only if $E$ is $m$-pseudo-coherent for
all $m$ (by definition) and the same is true for $M^\bullet$
by More on Algebra, Lemma \ref{more-algebra-lemma-pseudo-coherent}.
\end{proof}
